\documentclass[11pt]{article}
\usepackage[margin=1in]{geometry}
\usepackage{amsmath,amssymb,amsthm,mathtools}
\usepackage[T1]{fontenc}
\usepackage{lmodern,microtype}
\usepackage[colorlinks=true,linkcolor=blue,citecolor=blue,urlcolor=blue]{hyperref}
\usepackage{enumitem}
\usepackage{longtable}
\usepackage{xurl}
\allowdisplaybreaks
\setlength{\emergencystretch}{3em}
\newtheorem{theorem}{Theorem}[section]
\newtheorem{lemma}[theorem]{Lemma}
\newtheorem{proposition}[theorem]{Proposition}
\newtheorem{corollary}[theorem]{Corollary}
\theoremstyle{definition}\newtheorem{definition}[theorem]{Definition}
\theoremstyle{remark}\newtheorem{remark}[theorem]{Remark}
\newcommand{\R}{\mathbb R}
\newcommand{\E}{\mathbb E}
\newcommand{\Prob}{\mathbb P}
\newcommand{\1}{\mathbf 1}
\newcommand{\ip}[2]{\langle #1,#2\rangle}
\newcommand{\norm}[1]{\lVert #1\rVert}
\newcommand{\abs}[1]{\lvert #1\rvert}
\newcommand{\HS}{\mathrm F}
\newcommand{\Cov}{\operatorname{Cov}}
\newcommand{\Var}{\operatorname{Var}}
\newcommand{\tr}{\operatorname{tr}}
\newcommand{\dist}{\operatorname{dist}}
\newcommand{\Vol}{\operatorname{Vol}}
\newcommand{\diag}{\operatorname{diag}}
\newcommand{\supp}{\operatorname{supp}}
\newcommand{\im}{\operatorname{im}}
\newcommand{\Id}{I}
\newcommand{\cP}{\mathcal P}
\newcommand{\cO}{\mathcal O}
\newcommand{\gam}{\gamma}
\newcommand{\dd}{\,d}
\newcommand{\eps}{\varepsilon}
\title{Gaussian Covariance Concavity and\\ Simplex Optimality near Independence}
\author{}
\date{Local review manuscript --- 31 October 2026}
\begin{document}
\maketitle
\begin{abstract}
For every probability law $\nu$ with finite second moment, we prove that its
maximal covariance with a centered Gaussian is a concave function of the
Gaussian covariance matrix. For a nondegenerate simplex target the Hessian
is negative definite at every positive definite covariance. The proof uses
Gaussian integration on the strata of a translated simplex normal fan and an
exact factorization through the simplex boundary maps. A weighted choice of
cycle coordinates cancels the extra boundary term. Singular covariance lifting
and approximation extend the result to arbitrary targets.

For every fixed number of cells, we deduce the balanced Gaussian first-moment
inequality, its equality cases, and global centroid coercivity. Complete
transfer arguments then give regular-simplex optimality on dimension-independent
intervals adjacent to zero: a balanced maximum for positive correlation, and
a minimum over pairs with arbitrary matching masses for negative correlation.
The latter includes the single-partition problem without prescribed masses.
The intervals depend on the number of cells. Full-correlation ranges remain
open; no sharp multiway-cut hardness result is asserted.
\end{abstract}
\noindent\textit{Review status.}
This is a local manuscript for mathematical review. Its source lineage and
internal checks are recorded in Appendix~\ref{app:lineage}; no external review,
publication, or priority clearance is claimed.
\tableofcontents
\section{Statements and conventions}
Write $\gamma_N$ for standard Gaussian measure on $\R^N$ and $\cP_2(\R^N)$
for probability laws with finite second moment. For $\nu\in\cP_2(\R^N)$ set
\[
 \mathcal F_\nu(C)=\sup\{\E[X\cdot Y]:X\sim N(0,C),\ Y\sim\nu\},
 \qquad C\geq0.
\]
The supremum is over all couplings. Covariance matrices are varied in their
ordinary linear coordinates.

\begin{theorem}[Covariance concavity]\label{thm:covariance}
For every finite $N$ and $\nu\in\cP_2(\R^N)$, the function
$C\mapsto\mathcal F_\nu(C)$ is concave on the positive semidefinite cone.
Equivalently $C\mapsto W_2^2(N(0,C),\nu)$ is convex.
If $\nu$ has $d+1$ affinely independent sites in $\R^d$, each of positive
mass, then $\mathcal F_\nu$ is smooth on the positive definite cone and
its Hessian is negative definite there.
\end{theorem}

Fix $m\geq3$, put $d=m-1$, and identify
$H=\{x\in\R^m:\sum_i x_i=0\}$ isometrically with $\R^d$.
Let
\[
 v_i=e_i-\frac1m\1,\quad
 \mu_m=\E\max_{1\leq i\leq m}Z_i,\quad
 \alpha=\frac{\mu_m}{m-1},\quad
 c_i=\alpha v_i,\quad
 w=\frac{\mu_m^2}{m-1},
\]
where the $Z_i$ are independent standard normals. The regular partition is
$g_i(x)=\1_{\{v_i\cdot x=\max_j v_j\cdot x\}}$, with arbitrary choices on
null tie hyperplanes. Symmetry gives $\E[Xg_i(X)]=c_i$:
the centroids are multiples of $v_i$ and
$\sum_i v_i\cdot\E[Xg_i(X)]=\mu_m$ determines the multiplier.
In larger dimensions the partition ignores the unused coordinates.
With the noise notation below, every benchmark $S_\rho(g)=B_\rho(g,g)$ is
evaluated in dimension $d=m-1$ and is unchanged by adding unused Gaussian
coordinates. When $N<d$, the comparison is with this $d$-dimensional benchmark,
not with an assignment $g$ in $\R^N$.

A \emph{relaxed assignment} is a measurable map into
$\Delta_m=\{t\in[0,1]^m:\sum_i t_i=1\}$. It is balanced if
$\E f_i=1/m$. Write $b_i(f)=\E[Xf_i(X)]$ and
$W(f)=\sum_i\norm{b_i(f)}^2$. Let $\cO_m$ denote the ordered regular
centroid orbit under orthogonal maps. Label permutations are already
implemented by orthogonal symmetries of the simplex.

\begin{theorem}[Balanced endpoint and coercivity]\label{thm:endpoint}
For every balanced relaxed assignment in any finite dimension,
$W(f)\leq w$. If $N\geq d$, equality holds exactly for regular simplex
partitions, up to orthogonal maps, labels, and null changes. There is a finite
$K_m\geq1$, depending only on $m$, such that for all $N\geq d$,
\begin{equation}\label{eq:Hm}
 \dist_{\HS}(b(f),\cO_m)^2\leq K_m\bigl(w-W(f)\bigr).
\end{equation}
For $N<d$ the endpoint benchmark is strict.
\end{theorem}

For $|\rho|<1$ let
$T_\rho u(x)=\E u(\rho x+\sqrt{1-\rho^2}Z)$ and define
\[
 B_\rho(f,h)=\sum_i\ip{f_i}{T_\rho h_i}_{L^2(\gamma_N)},
 \qquad S_\rho(f)=B_\rho(f,f).
\]

\begin{theorem}[Near independence]\label{thm:noise}
For each fixed $m\geq3$ there are numbers $\rho_m^+>0$ and $r_m^->0$,
independent of ambient dimension, with the following properties in every
finite ambient dimension $N$.
\begin{enumerate}[label=(\roman*)]
\item If $0<\rho\leq\rho_m^+$ and $f$ is balanced, then
$S_\rho(f)\leq S_\rho(g)$.
\item If $0<r\leq r_m^-$ and $\E f=\E h=p$ for an arbitrary common
probability vector $p$, then $B_{-r}(f,h)\geq B_{-r}(g,g)$.
\end{enumerate}
For $N\geq d$ equality in (i) has the regular-simplex classification.
Equality in (ii) forces $p=(1/m,\ldots,1/m)$ and $f=h$ equal to one common
regular simplex partition. For $N<d$ both benchmarks are strict.
All statements include relaxed assignments.
\end{theorem}

In particular the negative single-assignment minimum has no imposed mass
constraint. This is not a theorem optimizing a prescribed unequal mass vector.
The two intervals can differ and may shrink with $m$; no numerical bound
uniform in $m$ is claimed. The covariance proof and endpoint argument are
given in the main text. Appendices~\ref{app:common}--\ref{app:negative} give
the complete fixed-dimensional transfers. Section~\ref{sec:dimension} states
the precise external dimension reductions and completes the ambient argument.

\paragraph{Related problems.}
Heilman \cite[Section 8, Conjecture 3]{H2014} asks whether balanced
four-cell first-moment maximizers in $\R^3$ are simplicial conical partitions,
and identifies quantitative stability as a route to larger-cell
near-zero noise results. Theorem~\ref{thm:endpoint} supplies this structure
and stability, and extends them to every $m\geq3$.
The introduction following Theorem 1.3 of Heilman's 23 October 2026
preprint \cite{H2026} still describes the four-cell and higher
Standard Simplex cases as open. That discussion distinguishes the
unrestricted Propeller problem, which does not prescribe cell masses,
from the Standard Simplex problem. It also distinguishes Gaussian
mean-width comparisons with prescribed individual variances.
Our endpoint prescribes cell masses and permits unequal centroid lengths;
it does not assume those variance constraints.
These comparisons locate the related problems and do not constitute an
exhaustive priority assessment. The full-correlation ranges remain open.


\section{Covariance concavity}
\subsection{The simplex dual and its strata}
First let $\nu=\sum_{i=0}^d p_i\delta_{y_i}$, where the sites are affinely
independent and $p_i>0$. For $C>0$,
\begin{equation}\label{eq:dual}
 \mathcal F_\nu(C)=\min_h
 \left\{\E\max_i(y_i\cdot X+h_i)-\sum_i p_i h_i\right\},
 \qquad X\sim N(0,C).
\end{equation}
Indeed the maximum bounds every coupling. The objective is at least
$\max h-\sum p_i h_i\geq(\min p_i)(\max h-\min h)$, so a minimum exists
modulo common offsets. Ties have Gaussian measure zero. Differentiation in
$h_i$ makes the winning masses equal to $p_i$, and that winning assignment
attains the bound.

Every offset class has a unique representation $h_i=c-y_i\cdot a$.
The cells are therefore the target simplex's normal fan translated by $a$.
For each nonempty vertex set $S$, let $\Gamma_S$ be the locus where all
labels in $S$ tie for the maximum. It is a simplicial cone in an affine
space of dimension $d+1-|S|$, based at $a$, with boundary facets
$\Gamma_{S\cup\{k\}}$, $k\notin S$. At $C=I$ define
\[
 \mu_S=\int_{\Gamma_S}\gamma_d(x)\dd\mathcal H^{d+1-|S|}(x),\qquad
 m_S=\frac1{\mu_S}\int_{\Gamma_S}x\gamma_d(x)\dd\mathcal H^{d+1-|S|}(x),
\]
and let $P_S$ be the orthogonal projection onto its tangent space.
All these masses are positive and all polynomial moments are finite.
For the apex, $\mu_{\{0,\ldots,d\}}=\gamma_d(a)$ and its mean is $a$.

The offset objective is the Gaussian convolution of the fixed support
function $\max_i y_i\cdot x$, evaluated at $-a$, plus
$(\sum_i p_i y_i)\cdot a$. Linear growth of the support function gives smooth
dependence on $a,C>0$ by differentiation of the Gaussian density. The offset
Hessian is the complete graph Laplacian with positive facet weights
$\mu_{\{i,j\}}/|y_i-y_j|$. It is invertible modulo constants. The implicit
function theorem thus gives smooth minimizing offsets and a well-defined
optimized Hessian.

\subsection{Intrinsic Gaussian integration}
For $i\in S$ and $k\notin S$, the inward normal within the stratum is
\[
 n_{S,k}=\frac{P_S(y_i-y_k)}{|P_S(y_i-y_k)|}.
\]
Intrinsic integration by parts yields, for positive-dimensional strata,
\begin{equation}\label{eq:stratum-mean}
 \mu_S P_Sm_S=\sum_{k\notin S}\mu_{S\cup\{k\}}n_{S,k}.
\end{equation}
For an edge stratum $e=\{i,j\}$ a second integration gives
\begin{equation}\label{eq:stratum-cov}
 \mu_e(P_e-\Cov_e X)
 =\sum_{k\notin e}\mu_{e\cup\{k\}}
       n_{e,k}(m_e-m_{e\cup\{k\}})^T.
\end{equation}
To verify the signs, the intrinsic Gaussian derivative is $-P_Sx\gamma_d$,
whereas the outward boundary normal is $-n_{S,k}$. In the second identity,
integrate the derivative of $((x-m_e)\cdot u)\gamma_d$ in a tangent direction
$v$. The two negative boundary signs produce the displayed product.
Both means lie in the same affine facet, so their difference is tangent.
Truncation by intrinsic balls produces only polynomial Gaussian tail terms
on the artificial sphere, which vanish as the radius tends to infinity.
Lower-dimensional intersections have zero boundary measure.

We use the classical convex-support Brascamp--Lieb consequence
\begin{equation}\label{eq:BL}
 \Cov_e X\leq P_e.
\end{equation}
For a direct applicable formulation, Bobkov--Ledoux
\cite[Proposition 2.1]{BL2000} allows a smooth strictly convex potential on
an open convex support. In orthonormal intrinsic coordinates the relative
interior of the facet is such a support, and its Gaussian potential has
Hessian $I$. Apply that proposition to linear functions; the facet boundary
has zero intrinsic measure. Neither strictness nor invertibility of the
covariance deficit is needed.

\subsection{Facet matrices}
Orient each edge $e=(i,j)$ with $i<j$ and put
\[
 d_e=y_i-y_j,\qquad \ell_e=|d_e|,\qquad w_e=\mu_e/\ell_e.
\]
For every target triangle $g$ let $A_g>0$ be its Euclidean area and set
$D_g=2A_g/\mu_g$. Let $l_k$ be the global affine barycentric coordinate
of vertex $k$ and write $v_k=-\nabla l_k$. Then
\begin{equation}\label{eq:barycentric}
 \sum_kv_k=0,\qquad v_k\cdot(y_i-y_j)=\delta_{kj}-\delta_{ki}.
\end{equation}
For $k\notin e$, $v_k$ is tangent to $\Gamma_e$. The vectors
$P_e(y_i-y_k)$ form a basis of that tangent space, with dual rows $v_k^T$.

Index the columns of $N_e,T_e$ by triangles $g=e\cup\{k\}$ and define
\[
 N_e[:,g]=\frac{\mu_g}{\ell_e}n_{e,k},\qquad
 T_e[:,g]=m_e-m_g,\qquad E_e=w_e(P_e-\Cov_e X).
\]
All matrices and inverses in this paragraph act on the facet tangent space.
The triangle height formula gives
\begin{equation}\label{eq:Ninverse}
 N_e[:,g]=\frac{P_e(y_i-y_k)}{D_g},
 \qquad \operatorname{row}_g(N_e^{-1})=D_gv_k^T.
\end{equation}
Equations~\eqref{eq:stratum-cov}--\eqref{eq:BL} imply
\begin{equation}\label{eq:Blocal}
 E_e=N_eT_e^T=T_eN_e^T\geq0,\quad
 T_e=E_eN_e^{-T},\quad
 B_e=N_e^{-1}E_eN_e^{-T}=N_e^{-1}T_e\geq0.
\end{equation}
There is no inverse of $E_e$ here. Equation~\eqref{eq:stratum-mean} also gives
$w_eP_em_e=N_e\1$, hence the exact entries
\begin{equation}\label{eq:Bentries}
 B_e[g,g']=\frac1{w_e}-D_gv_k\cdot m_{g'},\qquad g=e\cup\{k\}.
\end{equation}

\subsection{The chain factorization}
Let $\partial_2$ and $\partial_3$ be the usual boundary maps from oriented
triangles to edges and from tetrahedra to triangles, on the complete simplex.
All vertices in a simplex are ordered increasingly; for example
$\partial_2[i,j,k]=[j,k]-[i,k]+[i,j]$.
Coning chains to one fixed vertex proves that the graph cycle space is
$\im\partial_2$ and $\ker\partial_2=\im\partial_3$.
In dimension two $\partial_3$ is absent.

Put $W=\diag(w_e)$ and
$A=\partial_2^TW^{-1}\partial_2$. For a triangle coefficient vector
$\lambda$, let $\eta_{e,g}=(\partial_2)_{e,g}\lambda_g$ on triangles
incident to $e$, and define $B$ by
$\lambda^TB\lambda=\sum_e\eta_e^TB_e\eta_e$.
For a target tetrahedron $h$, let $V_h$ be its intrinsic three-dimensional
volume, and set $c_h=\mu_h/(6V_h)$. Also write $D=\diag(D_g)$.

\begin{lemma}[Exact boundary factorization]\label{lem:factor}
\begin{equation}\label{eq:factor}
 A-B=-D\partial_3\diag(c_h)\partial_3^TD.
\end{equation}
\end{lemma}
\begin{proof}
Suppose distinct triangles $g,g'$ share edge $e$, and let $k=g\setminus e$.
Because $k\notin g'$, the vector $v_k$ is tangent to $\Gamma_{g'}$.
Test~\eqref{eq:stratum-mean} there against $v_k$.
By~\eqref{eq:barycentric}, all boundary normals vanish in this pairing except
the one adjoining $k$. Its pairing is the reciprocal target height
$A_{g'}/(3V_h)$, where $h=g\cup g'$. Thus
\[
 v_k\cdot m_{g'}=\frac{\mu_hA_{g'}}{3V_h\mu_{g'}}.
\]
The off-diagonal entry of $A-B$ is therefore
$(\partial_2)_{e,g}(\partial_2)_{e,g'}D_gD_{g'}c_h$.
The identity $\partial_2\partial_3=0$, on the shared edge of this tetrahedron,
gives
\[
 (\partial_2)_{e,g}(\partial_2)_{e,g'}
 =-(\partial_3)_{g,h}(\partial_3)_{g',h}.
\]
This proves the claimed off-diagonal entries. Triangles sharing no edge
interact in neither matrix.

For the diagonal, summing~\eqref{eq:Bentries} over the three edges of $g$
and using~\eqref{eq:barycentric} gives
\[
 (A-B)_{g,g}
 =D_g\sum_{k\in g}v_k\cdot m_g
 =-D_g\sum_{l\notin g}v_l\cdot m_g
 =-D_g^2\sum_{h\supset g}c_h.
\]
The last equality uses the same stratum mean identity and reciprocal height.
If $d=2$, the sum over outside vertices is empty and the difference is zero,
also when the triangle stratum is the apex. These are all entries of
\eqref{eq:factor}. The volumes in this argument are intrinsic to the target
tetrahedra, so there is no factor depending on the ambient dimension.
\end{proof}

\subsection{An admissible gauge and the cycle estimate}
For a fixed edge cycle $j$, choose $\partial_2\lambda=j$.
Its representatives form $\lambda+\im\partial_3$. Minimize
$\lambda^TD\lambda$ on this closed affine space. Since $D>0$, the minimizing
representative exists and satisfies
\begin{equation}\label{eq:gauge}
 \partial_3^TD\lambda=0.
\end{equation}
This uses no injectivity of $\partial_3$. It changes only cycle coordinates.
By~\eqref{eq:factor},
\begin{equation}\label{eq:gauge-equality}
 \lambda^TB\lambda=\lambda^TA\lambda=j^TW^{-1}j.
\end{equation}

Give the cells arbitrary vectors $z_i$ and write
$z_e=z_i-z_j$, $a_e=z_e\cdot m_e$.
Signed edge differences telescope around each triangle, so
\[
 j\cdot a=\sum_e z_e\cdot T_e\eta_e.
\]
By~\eqref{eq:Blocal}, square-root Cauchy--Schwarz in the PSD matrices $E_e$,
followed by Cauchy--Schwarz across edges, gives
\[
 (j\cdot a)^2\leq
 \left(\sum_e z_e^TE_ez_e\right)j^TW^{-1}j.
\]
Singular deficits cause no exception. Weighted least-squares duality on the
edge graph now yields
\begin{equation}\label{eq:cycle-bound}
 \min_b\sum_e w_e(a_e+b_i-b_j)^2
 \leq\sum_e w_e\bigl(|P_ez_e|^2-\Var_e(z_e\cdot X)\bigr).
\end{equation}
For clarity, at the least-squares minimizer the weighted residual is a cycle.
Cauchy--Schwarz gives the dual supremum
$\sup_{j\ne0}(j\cdot a)^2/(j^TW^{-1}j)$, and that residual attains it
unless it is zero. Thus~\eqref{eq:cycle-bound} uses a single global offset
vector $b$, not offsets chosen independently on facets.

\subsection{The exact covariance Hessian}
Take $C(t)=I+tH$ with $H=H^T$. The first two derivatives of its square root
at zero are $H/2$ and $-H^2/4$. Differentiating the affine maximum in
\eqref{eq:dual} and eliminating the offset derivative gives
\begin{equation}\label{eq:cov-Hessian}
 4D^2\mathcal F_\nu(I)[H,H]
 =\min_b\sum_e w_e\E\bigl[(X\cdot Hd_e+b_i-b_j)^2\mid\Gamma_e\bigr]
       -\sum_e w_e|Hd_e|^2.
\end{equation}
Here conditional expectation means the normalized intrinsic facet measure.
To detail the second variation, the score velocity difference across a facet
is $X\cdot Hd_e/2+h_i'-h_j'$. Differentiating the first derivative over moving
cells contributes its square divided by $\ell_e$, integrated over the facet.
Each unoriented facet is counted once. Gaussian integration by parts for the
winning vector gives
$\E[XY(X)^T]=\sum_e w_ed_ed_e^T$, so the interior acceleration is
$-\frac14\sum_e w_e|Hd_e|^2$. Put $b=2h'$.
The $h''$ term vanishes because the current masses equal $p$.
The positive offset Hessian makes the displayed minimum exactly the Schur
complement. Polyhedral truncation and Gaussian tails justify the boundary
differentiation; multiple ties have zero facet measure, and the smoothness
already proved identifies this calculation with the ordinary Hessian.

Apply~\eqref{eq:cycle-bound} with $z_i=Hy_i$. The facet variances cancel,
leaving
\begin{equation}\label{eq:strict-Hessian}
 4D^2\mathcal F_\nu(I)[H,H]
 \leq-\sum_e w_e\frac{(d_e^THd_e)^2}{\ell_e^2}.
\end{equation}
The edge dyads of an affinely independent simplex span the symmetric matrices:
differences $b_i=y_i-y_0$ form a basis, and the quadratic forms at $b_i$
and $b_i-b_j$ recover every bilinear entry in that basis.
Thus the right side is strictly negative for $H\ne0$.
At any $C_0>0$, transform the target sites by $C_0^{1/2}$ and the covariance
direction $K$ by $H=C_0^{-1/2}KC_0^{-1/2}$. This reduces to the calculation
at $I$, without a commutativity assumption. Integration along segments proves
concavity on the positive definite cone. For $d=1$ the value is a positive
constant times the square root of variance, giving the same assertions.

\subsection{Singular covariance, finite lifting, and limits}
For every fixed $\nu\in\cP_2$,
\begin{equation}\label{eq:C-continuity}
 |\mathcal F_\nu(C)-\mathcal F_\nu(C')|
 \leq(\E|Y|^2)^{1/2}\norm{C^{1/2}-(C')^{1/2}}_{\HS}.
\end{equation}
Indeed any coupling of $C^{1/2}G$ with $Y$ lifts to one with standard $G$,
using its Gaussian conditional law even when $C$ is singular. Reuse $G,Y$
for $C'$ and apply Cauchy--Schwarz; then reverse the argument.
Adding $\eps I$ and taking limits extends simplex concavity to the PSD cone.

For an arbitrary finite target $\sum_{i=1}^m p_i\delta_{y_i}$ in $\R^N$,
lift its sites to $(y_i,\eps e_i)$ in $\R^{N+m}$, for any fixed $\eps>0$.
They are affinely independent. A Gaussian with covariance $\diag(C,0)$
has exactly the old cost, so the two optimal values coincide. Subtract a
common lifted site; its expected contribution is zero. Project onto the
$(m-1)$-dimensional span $S$ of lifted differences. The projected source has
covariance $P_S\diag(C,0)P_S$, linear in $C$, and its cost is unchanged.
Every projected coupling lifts by the source conditional law. Simplex PSD
concavity therefore pulls back to the original finite-target functional.
A one-atom target has value zero.

Finally, gluing a coupling of two targets to a Gaussian-target coupling gives
\begin{equation}\label{eq:target-continuity}
 |\mathcal F_\nu(C)-\mathcal F_{\nu'}(C)|
 \leq\sqrt{\tr C}\,W_2(\nu,\nu').
\end{equation}
Finite laws are dense in $\cP_2$, by truncation and quantization. Pass to the
limit in the concavity inequality. The identity
$W_2^2(N(0,C),\nu)=\tr C+\int|y|^2\dd\nu-2\mathcal F_\nu(C)$ proves the
equivalent assertion in Theorem~\ref{thm:covariance}. Neither lifting nor
limits assert strictness for general targets.

\begin{corollary}\label{cor:square}
The function $C\mapsto\mathcal F_\nu(C)^2$ is concave on the PSD cone.
\end{corollary}
\begin{proof}
Independence gives $\mathcal F_\nu\geq0$, and scaling the Gaussian gives
$\mathcal F_\nu(tC)=\sqrt t\,\mathcal F_\nu(C)$.
For $0<a<1$, concavity applied to $C/a$ and $D/(1-a)$ yields
\[
 \mathcal F_\nu(C+D)\geq
 \sqrt a\,\mathcal F_\nu(C)+\sqrt{1-a}\,\mathcal F_\nu(D).
\]
Optimizing $a$, with endpoints obtained by limits, proves superadditivity
of $\mathcal F_\nu^2$. Its degree-one homogeneity then proves concavity.
\end{proof}


\section{The balanced endpoint and a global coercivity estimate}
Let $P=I-\1\1^T/m$, $K_*=mP$, and $L_0=m(m-1)$. On zero-sum score
covariances define
\[
 J_m(K)=\inf_{\sum_i h_i=0}\E\max_i(X_i+h_i),\qquad
 X\sim N(0,K),\quad K\geq0,\ K\1=0.
\]
This is the maximal covariance functional for the uniform regular target
$v_1,\ldots,v_m$ on the zero-sum space. The dual remains valid for singular
$K$ by continuity. Indeed the difference between two expected maxima,
coupled through one standard Gaussian and with the same offsets, is at most
$\norm{K^{1/2}-(K')^{1/2}}_{\HS}$, uniformly in those offsets. The same bound
therefore controls their infima. Approximate on the zero-sum space by positive
definite covariances and also use~\eqref{eq:C-continuity}.
The assignment inequality below follows directly from the pointwise maximum,
irrespective of singular dual attainment.

On $\tr K=L_0$, permutation averaging gives $K_*$: invariance forces constant
diagonal and off-diagonal entries, fixed by the trace and zero row sums.
Theorem~\ref{thm:covariance} and permutation invariance therefore imply
\begin{equation}\label{eq:score-maximum}
 J_m(K)\leq J_m(K_*)=\sqrt m\,\mu_m=\sqrt{L_0w}.
\end{equation}
At $K_*$ the winner assignment is balanced by symmetry, so zero offsets
attain this value.

The strict Hessian in~\eqref{eq:strict-Hessian}, continuity, and permutation
symmetry give constants $c_0,\delta_0>0$ such that
\[
 J_m(K_*)-J_m(K)\geq c_0\norm{K-K_*}_{\HS}^2
 \quad\hbox{when }\norm{K-K_*}_{\HS}\leq\delta_0,\quad \tr K=L_0.
\]
Indeed the gradient at $K_*$ is scalar on the zero-sum space, so the first
derivative vanishes on trace-zero directions; integrate the negative Hessian
in a small neighborhood. The trace-normalized PSD domain is compact. Choose
one $\theta\in(0,1)$ so that
$K_\theta=(1-\theta)K_*+\theta K$ lies in that neighborhood for every $K$
in the domain. Concavity gives
\[
 \theta[J_m(K_*)-J_m(K)]
 \geq J_m(K_*)-J_m(K_\theta)
 \geq c_0\theta^2\norm{K-K_*}_{\HS}^2.
\]
Thus, with $\kappa_m=c_0\theta>0$,
\begin{equation}\label{eq:global-Gram-gap}
 J_m(K_*)-J_m(K)\geq\kappa_m\norm{K-K_*}_{\HS}^2
\end{equation}
on the entire normalized PSD domain, including its boundary.

For balanced relaxed $f$ with $W=\sum_i|b_i|^2>0$, normalize
$B_i=\sqrt{L_0/W}\,b_i$ and let $K=BB^T$. Since $\sum_i b_i=0$,
$K\1=0$ and $\tr K=L_0$. The same assignment is an admissible balanced
score coupling, hence
\[
 \sqrt{L_0W}=\E\sum_i f_i(X)B_i\cdot X
 \leq J_m(K)\leq\sqrt{L_0w}.
\]
This proves $W\leq w$; $W=0$ is immediate.
If $W=w$, \eqref{eq:global-Gram-gap} forces $K=K_*$. Equality in the
pointwise maximum inequality then forces $f$ to select the unique regular
winning score almost everywhere. This also proves the relaxed equality
classification. In $N<d$, pad the source by independent Gaussian coordinates;
the centroid Gram has rank at most $N$ and cannot equal a regular rank-$d$
Gram, proving strictness.

For quantitative centroid alignment, put $\Delta=w-W$. The difference of
square roots and~\eqref{eq:global-Gram-gap} imply
\[
 \norm{K-K_*}_{\HS}^2
 \leq\frac{\sqrt{L_0}}{\kappa_m\sqrt w}\Delta.
\]
For $N\geq d$, orthogonal Procrustes alignment with any regular row matrix
$B_*$ of Gram $K_*$ gives
\begin{align*}
 \min_{O\in O(N)}\norm{B-B_*O}_{\HS}^2
 &=\sum_{j=1}^{m-1}(\sqrt{\lambda_j(K)}-\sqrt m)^2\\
 &\leq \frac1m\norm{K-K_*}_{\HS}^2.
\end{align*}
Rescale the two matrices by $\sqrt{W/L_0}$ and $\sqrt{w/L_0}$ respectively.
The two-square inequality, $W\leq w$, and
$(\sqrt W-\sqrt w)^2\leq\Delta$ show that
\begin{equation}\label{eq:K-explicit}
 \dist_{\HS}(b,\cO_m)^2
 \leq\left(2+\frac{2\sqrt w}{m\kappa_m\sqrt{L_0}}\right)\Delta.
\end{equation}
The case $W=0$ is immediate with this coefficient. This proves
Theorem~\ref{thm:endpoint}, with $K_m$ equal to the displayed coefficient or
any larger number. It is finite and independent of $N$; no explicit
numerical value or uniform bound as $m\to\infty$ is asserted.


\section{Ambient dimensions and discrete consequences}\label{sec:dimension}
The proofs in Appendices~\ref{app:positive} and~\ref{app:negative} take place
in $d=m-1$. We use the following two external results with different feasible
classes.

Heilman--Tarter \cite[Theorem 1.9]{HT2021} states that, for positive
correlation and prescribed positive cell masses, every maximizing Gaussian
partition in dimension $N\geq m-1$ depends, after a common orthogonal change
of coordinates, on at most $m-1$ coordinates.
Heilman \cite[Problem 7.1 and Theorem 7.9]{H2022} instead treats minimizers of
a positive-kernel bilinear objective over two partitions whose corresponding
masses agree but are \emph{free}. Every such minimizer has a common
$(m-1)$-dimensional reduction. This latter result has no hyperstability
hypothesis and is not the distinct theorem for two separately prescribed
mass vectors.

\subsection{Existence and the positive problem}
For finite $N$, the simplex-valued feasible set with prescribed masses is
weakly compact in $L^2(\gamma_N)^m$. The Ornstein--Uhlenbeck operator is compact:
each Hermite degree has finite multiplicity and $|\rho|^k\to0$.
Thus the positive quadratic objective is weakly continuous on this bounded
set and attains its maximum.

Every maximizer is an indicator partition. Otherwise two fixed labels are
both uniformly fractional on a positive-measure set. Choose two disjoint
subsets of equal positive measure there. Opposite label perturbations on
those subsets preserve all masses and produce two distinct feasible
assignments whose average is the original one. Strict positivity of $T_\rho$
for $\rho>0$, visible from its strictly positive Hermite eigenvalues, makes
the average of their objectives strictly larger. This contradicts maximality.
Theorem~1.9 of \cite{HT2021} therefore applies to every relaxed maximizer.
Appendix~\ref{app:positive} gives its value and its exact equality form.

\subsection{Free matching masses and the negative problem}
The set of pairs with free common masses is weakly compact, and compactness
of $T_{-r}$ makes the bilinear objective jointly weakly continuous on it.
A minimizing pair exists. Hold its attained mass vector $p$ fixed.
With $h$ fixed, the minimizing face for the continuous linear functional in
$f$ has an extreme point. Every extreme point of this fixed-mass
simplex-assignment set is an indicator, by the same equal-measure opposite
perturbation construction. Choose such an $f_1$, and then with $f_1$ fixed
choose an indicator minimizer $h_1$. The pair has the same global minimum.
This sequential argument does not require an incorrect general principle
about extreme points of bilinear optimization.

The competitive bounds~\eqref{eq:competitive-negative} hold in every finite
ambient dimension, and~\eqref{eq:negative-range} makes all masses positive.
Reflect the second partition: $B_{-r}(f_1,h_1)=B_r(f_1,h_1(-\,\cdot))$.
The free-common-mass theorem \cite[Theorem 7.9]{H2022} now reduces this
indicator minimizer to a common subspace of dimension at most $d$.
Reflection back preserves that subspace. Appendix~\ref{app:negative} forces
$f_1=h_1=g$ in one common frame.

For the original fractional pair, the intermediate pair $(g,h)$ still has
the minimum value. Its excess over $(g,g)$ is exactly $D_-(h)$.
Strict reference ordering forces $h=g$. The original pair $(f,g)$ then
has excess $D_-(f)$, so $f=g$. This recovers equality for the original
fractional pair, not only for an indicator representative.
In either sign, an assignment in $N<d$ can be embedded in dimension $d$;
equality would require a regular centroid Gram of rank $d$, impossible for
the embedded rank at most $N$. This completes Theorem~\ref{thm:noise}.

\subsection{The precise low-influence corollary}
On $[m]^n$ use uniform product measure and the noise channel which, in each
coordinate, keeps a symbol with probability
$(1+(m-1)\rho)/m$ and changes it to each particular different symbol with
probability $(1-\rho)/m$. It exists for $-1/(m-1)\leq\rho\leq1$.
For $f:[m]^n\to\Delta_m$ define
\[
 \operatorname{Inf}_k(f)=\sum_i\E\Var(f_i\mid\hbox{all coordinates except }k).
\]
Write $\operatorname{Stab}_\rho(f)=\sum_i\E[f_i(X)f_i(Y)]$ for this channel,
and $p_m(\rho)=S_\rho(g)$.

\begin{corollary}\label{cor:discrete}
Fix $m\geq3$, an admissible nonzero $\rho$ in the corresponding interval of
Theorem~\ref{thm:noise}, and $\eps>0$. There is $\tau>0$, depending on
$m,\rho,\eps$, such that for every $n$ and every $f$ with
$\max_k\operatorname{Inf}_k(f)\leq\tau$,
\[
 \operatorname{Stab}_\rho(f)\leq p_m(\rho)+\eps
 \quad(\rho>0,\ \E f=\1/m),
\]
whereas
\[
 \operatorname{Stab}_\rho(f)\geq p_m(\rho)-\eps
 \quad(\rho<0,\ \text{with arbitrary }\E f).
\]
\end{corollary}
\begin{proof}
Isaksson--Mossel \cite[Theorem 7.1]{IM2009} transfers a low-influence
simplex-valued function to a Gaussian relaxed assignment with the same
expectation vector and approximates stability at the same fixed $\rho$.
Apply Theorem~\ref{thm:noise}, using its free-mass single-assignment consequence
for the negative case, and allocate the approximation error below $\eps$.
Lemma~7.3 of \cite{IM2009} identifies $p_m(\rho)$ with the limiting plurality
stability. This gives the stated additive comparison.
\end{proof}
There is no uniform-in-$\rho$ influence threshold or assertion of exact
finite-electorate optimality. Only the single-rule arbitrary-bias negative
corollary is asserted here; no additional discrete bilinear theorem is needed.
The near-zero Gaussian intervals do not give a full-range Standard Simplex
theorem or a sharp MAX-$m$-CUT hardness ratio.


\appendix
\section{Common estimates for the two transfers}\label{app:common}
Throughout the transfer appendices $d=m-1$, the reference $g$ and centroids
$c_i$ are as in Section~1, and $K=K_m\geq1$ satisfies~\eqref{eq:Hm}.
Set
\[
 a=|c_i-c_j|=\sqrt2\alpha,\qquad
 P_0=\binom m2,\qquad
 \beta=P_0\sqrt{2/\pi}/a,\qquad \sigma=(K+1)^{-1}.
\]
The first moment obeys $w\leq1$ by Parseval, so
$a^2=2w/(m-1)\leq1$.
All integrals in these appendices use $\gamma_d$ unless specified otherwise.

\subsection{The same nearest frame controls the mismatch}
For balanced $f$ choose a nearest regular centroid matrix $c$ and its partition
$g$. Put $\Delta=w-W(f)$, $q=\norm{b(f)-c}_{\HS}^2$,
$\eta=\frac12\sum_i\norm{f_i-g_i}_1$, and
$J=w-\sum_i c_i\cdot b_i(f)$.
Then
\begin{equation}\label{eq:linear-identity}
 2J=\Delta+q\leq(1+K)\Delta.
\end{equation}
If $\Gamma$ is the gap between the two largest regular linear scores, a union
bound on the $P_0$ Gaussian score differences of standard deviation $a$ gives
$\Prob(\Gamma\leq t)\leq\beta t$. With $p_{\rm wrong}=1-f_{\rm winner}$,
$J\geq\E[p_{\rm wrong}\Gamma]$ and $\E p_{\rm wrong}=\eta$.
Layer cake yields
\[
 J\geq\int_0^\infty(\eta-\beta t)_+\dd t
   =\frac{\eta^2}{2\beta},\qquad
 \eta^2\leq L\Delta,\quad L=\beta(1+K).
\]
Thus no second choice of alignment is needed.

\subsection{A saturated positive gap and a damped negative gap}
Let $\phi,\Phi$ be the standard normal density and distribution function, and
\[
 H_\delta(u)=\phi(u)\Phi(u+\delta)-\phi(u+\delta)\Phi(u).
\]
For $\delta>0$, this is positive: $\Phi(u)/\phi(u)$ is strictly increasing,
as is clear from its integral representation
$\int_0^\infty e^{uy-y^2/2}\dd y$.
In reference cell $i$, put $z_{ij}=(v_i-v_j)\cdot x/\sqrt2\geq0$.
For independent normal utilities with label $i$ of mean zero and label $j$
of mean $-\delta$, the winning-probability difference is
\[
 \E\left[\int_M^\infty H_\delta(u)\dd u\right],
\]
where $M$ is the maximum of the other utilities. Lowering their means
increases this difference.

For $G_i=T_rg_i$, subtract the mean of winner $i$.
The other means are nonpositive and
$\delta=\sqrt2\,r z_{ij}/\sqrt{1-r^2}$.
Raise all other means to zero; the event $M\leq0$ has probability
$2^{-(m-2)}$. For $0\leq u,\delta\leq1$,
$\partial_\delta H_\delta(u)\geq\phi(1)\phi(2)$.
Integrate over $0\leq u\leq1$; for larger $\delta$ use monotonicity when the
competitor mean is lowered. Consequently
\begin{equation}\label{eq:plus-gap}
 G_i-G_j\geq c_+\min(rz_{ij},1),\qquad
 c_+=\frac{e^{-5/2}}{2^{m-1}\pi}.
\end{equation}

For $G_i=T_{-r}g_i$, subtract label $j$'s mean instead.
Now $j$ has mean zero, $i$ has mean $-\delta$, and every other mean is at most
$U_x=\sqrt2\,r|x|/\sqrt{1-r^2}$. Raise those means to $U_x$.
Their maximum is at most $U_x$ with probability $2^{-(m-2)}$.
On $U_x\leq u\leq U_x+1$ and $0\leq\delta\leq1$,
\[
 \partial_\delta H_\delta(u)
 =\phi(u)\phi(u+\delta)+(u+\delta)\phi(u+\delta)\Phi(u)
 \geq\phi(U_x+1)\phi(U_x+2).
\]
The same integral argument and monotonicity for $\delta\geq1$ give a gap at
least $c_+e^{-U_x^2-3U_x}\min(\delta,1)$. Since
$3U_x\leq U_x^2+9/4$ and
$2U_x^2=4r^2|x|^2/(1-r^2)\leq|x|^2/32$ for $r\leq1/16$,
\begin{equation}\label{eq:minus-gap}
 G_j-G_i\geq c_-\min(rz_{ij},1)e^{-|x|^2/32},
 \qquad c_-=\frac{e^{-19/4}}{2^{m-1}\pi}.
\end{equation}
Thus the reference winner is the unique largest positive potential and the
unique smallest negative potential, outside null hyperplanes.

\subsection{Weighted moments and uniform tails}
Let
\[
 \psi(x)=e^{|x|^2/8},\qquad
 \omega_r(x)=\min\left(r\min_{i<j}|(v_i-v_j)\cdot x|/\sqrt2,\,1\right).
\]
Use $\widetilde\omega_r^+=\omega_r$ and
$\widetilde\omega_r^-=e^{-|x|^2/32}\omega_r$.
Define the finite constants
\begin{align*}
 B_+&=P_0\left[2^{(d-1)/2}\phi(0)e^{1/4}+2^{d/2-2}\right],\\
 B_-&=P_0\left[\phi(0)e^{5/16}(8/3)^{(d-1)/2}
                  +\tfrac14(8/3)^{d/2}\right].
\end{align*}
For $|h|\leq1$, $\lambda>0$, and $0<r\leq1$,
\begin{equation}\label{eq:Fenchel}
 \int|h|\psi\leq\lambda\int|h|\widetilde\omega_r^\pm
                      +\frac{B_\pm}{\lambda r}.
\end{equation}
Here are the constants explicitly. For the positive case dominate the
positive part by summing over the individual planes. Write $x=ue+z$.
On $|u|\leq1$, replace $\psi$ by
$A(z)=e^{1/8}e^{|z|^2/8}$ and $\phi(u)$ by $\phi(0)$, and integrate
$(A-\lambda r|u|)_+$ over the whole line. Its integral is
$A^2/(\lambda r)$. On $|u|>1$ the saturated plane weight is at least $r$,
and $(A-B)_+\leq A^2/(4B)$. Gaussian exponential moments give $B_+$.
For the negative case use the exact identity
\[
 (\psi-\lambda\widetilde\omega_r^-)_+
 =e^{-|x|^2/32}(e^{5|x|^2/32}-\lambda\omega_r)_+
 \leq(e^{5|x|^2/32}-\lambda\omega_r)_+.
\]
The same calculation, with this auxiliary exponential, gives $B_-$.
The moment controlled on the left of~\eqref{eq:Fenchel} remains $\psi$.

For any relaxed $k$, define the nonnegative losses
\[
 D_+(k)=-\sum_i\ip{k_i-g_i}{T_rg_i},\qquad
 D_-(k)= \sum_i\ip{k_i-g_i}{T_{-r}g_i}.
\]
Each vanishes only for $k=g$ by the strict gaps. Optimizing
\eqref{eq:Fenchel}, using $\sum_i|k_i-g_i|=2p_{\rm wrong}$ and
\eqref{eq:plus-gap} or~\eqref{eq:minus-gap}, gives
\begin{equation}\label{eq:weighted-moments}
 \sum_i\left(\int|k_i-g_i|\psi\right)^2\leq M_\pm D_\pm(k)/r,
 \qquad M_\pm=8B_\pm/c_\pm.
\end{equation}
In the relevant sign put $A_\pm^2=\max(1,m/4)M_\pm$.
If $z=\sqrt{D_\pm(k)/r}$ and $\eta=\frac12\sum_i\norm{k_i-g_i}_1$,
then $\eta\leq A_\pm z$.
For $R\geq2$, let $\eps_R=(1+R)e^{-R^2/8}$.
The pointwise bound
$(1+|x|)\1_{\{|x|>R\}}\leq\eps_R\psi(x)$ implies both
\begin{equation}\label{eq:tails}
 \frac12\sum_i\left|\E[(k_i-g_i)\1_{\{|X|>R\}}]\right|
 \leq\eps_RA_\pm z,\qquad
 \left(\sum_i\left|\E[X(k_i-g_i)\1_{\{|X|>R\}}]\right|^2\right)^{1/2}
 \leq\eps_RA_\pm z.
\end{equation}
These are separate bounds, obtained from Cauchy--Schwarz across labels.

\subsection{The quadratic Mehler error}
The Mehler kernel relative to $\gamma_d\otimes\gamma_d$ is
\[
 M_\rho(x,y)=(1-\rho^2)^{-d/2}
 \exp\!\left(\frac{2\rho x\cdot y-\rho^2(|x|^2+|y|^2)}
                   {2(1-\rho^2)}\right).
\]
For $|\rho|\leq1/16$, writing $\ell=\log M_\rho$ and
$L_{xy}=|x|^2+|y|^2$, direct differentiation gives
$|\ell'|\leq d+L_{xy}$, $|\ell''|\leq2d+2L_{xy}$ and
$M_\rho\leq2^d e^{L_{xy}/34}$.
For example the nondimensional parts of $\ell'$ and $\ell''$ are
\[
 \frac{(1+\rho^2)x\cdot y-\rho L_{xy}}{(1-\rho^2)^2},
 \qquad
 \frac{2\rho(3+\rho^2)x\cdot y-(1+3\rho^2)L_{xy}}{(1-\rho^2)^3};
\]
$|x\cdot y|\leq L_{xy}/2$ gives the stated bounds. The exponent bound follows
from $|\rho|/[2(1+|\rho|)]\leq1/34$.
Put $t_0=13/136$ and
\[
 C_d=2^{d-1}\left[d^2+2d+\frac{2d+2}{2t_0}+\frac1{t_0^2}\right].
\]
Taylor's theorem, $M''=M((\ell')^2+\ell'')$, and the maxima of
$L_{xy}e^{-t_0L_{xy}}$ and $L_{xy}^2e^{-t_0L_{xy}}$, using $e>2$, give
\begin{equation}\label{eq:Mehler}
 |M_\rho-1-\rho x\cdot y|
 \leq C_d\rho^2\psi(x)\psi(y).
\end{equation}
Thus a mixed remainder for deviations $f-g,h-g$ is bounded by
$C_dM_\pm r\sqrt{D_\pm(f)D_\pm(h)}$.

\subsection{Exact mass repair on fixed facet slabs}\label{subsec:repair}
Fix a reference frame, $R\geq3$, and a relaxed $k$ with mismatch $\eta$.
Let $p=\E k$, $U=\1/m$, $u=\frac12\norm{p-U}_1$.
Truncate $k$ to $g$ outside $B_R$, obtaining $k^R$, whose half-$\ell^1$
mass imbalance $v$ satisfies
\[
 v\leq u+\eps_RA_\pm z,\qquad v\leq\eta.
\]
Set
\[
 c_A=\Vol(B_{d-1})\,8^{-(d-1)}(2\pi)^{-d/2}e^{-2}>0.
\]
For donor $i$ and recipient $j$, use the facet disk centered at $v_i+v_j$,
of tangential radius $1/8$, and its donor-side normal slab of width
$\tau=(\eta+2v)/c_A$.
Assume $\eta\leq c_A/24$, so $\tau\leq1/8$.
At the anchor, the \emph{unscaled} $v_i$ and $v_j$ scores tie and exceed
all other scores by one. A displacement of size at most $\sqrt2/8$ reduces
these gaps by at most $1/4$. Distinct unoriented anchors are at least
$\sqrt2$ apart; opposite slabs at one anchor have disjoint interiors.
All oriented supports therefore lie in their correct cells, are mutually
disjoint up to null sets, and are contained in $B_2$.

Their Gaussian mass is at least $c_A\tau$.
Since the total missing donor mass in its reference cell is at most $\eta$,
each slab has available $k_i^R$ mass at least $c_A\tau-\eta=2v$.
Choose any excess-to-deficit flow of total mass $v$, and transfer each pair
mass proportionally to $k_i^R$ on its slab. Each pair amount is at most $v$.
Disjoint supports keep the assignment simplex-valued, and all masses are
repaired exactly. If $v=0$, make no transfers. Call the result $k'$.

Let $b'=b(k')$, $e_0=\norm{b'-b(k)}_{\HS}$, and
\[
 J'=w-\sum_i c_i\cdot b_i',\qquad
 J_{\rm in}=\E\!\left[\1_{B_R}\sum_i k_i(X)
                  (\max_j c_j\cdot X-c_i\cdot X)\right].
\]
The linear cost per unit transferred mass is at most $a\tau\leq\tau$.
Its total is at most $3\eta v/c_A$. Tail removal and the centroid cost of
transferring mass in $B_3$ consequently give
\begin{equation}\label{eq:repair-bounds}
 J'\leq J_{\rm in}+\frac3{c_A}\eta u+C_J\eps_Rz^2,\qquad
 e_0\leq6(u+\eps_RA_\pm z),\qquad C_J=3A_\pm^2/c_A.
\end{equation}
For the first inequality use $v\leq u+\eps_RA_\pm z$ and
$\eta\leq A_\pm z$; for the second use~\eqref{eq:tails} and
$3\sqrt2v$ as the Frobenius centroid cost.

Suppose now that $c$ is a nearest regular centroid matrix for $b(k)$,
and put $q=\norm{b(k)-c}_{\HS}^2$.
Distance to $\cO_m$ is one-Lipschitz. Applying~\eqref{eq:Hm} to balanced $k'$
gives
$2J'\geq(1+1/K)(\sqrt q-e_0)_+^2$.
Young's inequality with parameter $1/(2K)$, also covering $\sqrt q\leq e_0$,
therefore gives
\begin{equation}\label{eq:repair-coercivity}
 q\leq(2-\sigma)J'+(2K+1)e_0^2.
\end{equation}
The linear cost retains the original frame; no new alignment is introduced.

\subsection{Comparison on a fixed ball}
Write $H_i(y)=\E g_i(y+Z)$. Then $\nabla H_i(0)=c_i$ and
$\norm{\nabla^2H_i}_{\rm op}\leq\sqrt2$, by Gaussian differentiation:
in a unit direction the second derivative is
$\E[g_i(y+Z)((u\cdot Z)^2-1)]$, bounded by Cauchy--Schwarz.
Conicality gives $T_{\pm r}g_i(x)=H_i(\pm tx)$,
$t=r/\sqrt{1-r^2}$. On $B_R$ and for $r\leq1/16$,
\[
 \norm{\nabla(T_{\pm r}g_i-1/m\mp r c_i\cdot x)}
 \leq(1+3R)r^2.
\]
The orthogonal reflection swapping vertices $i,j$ fixes their bisector and
all other vertices, so the pair remainder vanishes on the bisector.
Integrating normally from that plane, inside $B_R$, gives the appropriate
signed potential gap at least
$r(c_i-c_j)\cdot x(1-\beta_Rr)$ in reference cell $i$, where
\begin{equation}\label{eq:compact}
 \beta_R=\frac{2(1+3R)}a,\qquad
 rJ_{\rm in}\leq\frac{D_\pm(k)}{1-\beta_Rr}
 \quad\hbox{if }\beta_Rr<1.
\end{equation}


\section{The positive-correlation transfer}\label{app:positive}
Work first in dimension $d=m-1$. Call balanced $f$ competitive if
$S_\rho(f)\geq S_\rho(g)$, and choose a nearest reference centroid frame.
Write $D=D_+(f)$, $z=\sqrt{D/\rho}$, and use the positive constants from
Appendix~\ref{app:common}. If $\Delta=w-W(f)=0$, the endpoint equality already
gives $f=g$, so assume otherwise.

Hermite expansion, positivity of its coefficients for $\rho>0$, and Parseval
give for a competitive assignment
\[
 \frac1m+\rho w\leq S_\rho(g)\leq S_\rho(f)
 \leq\frac1m+\rho W(f)+\rho^2(1-1/m).
\]
Hence $\Delta\leq(1-1/m)\rho\leq\rho$ and $\eta^2\leq L\rho$.
For $h=f-g$, the mean of each $h_i$ is zero. Equations~\eqref{eq:weighted-moments}
and~\eqref{eq:Mehler} give the exact comparison and its error bound
\begin{equation}\label{eq:positive-comparison}
 S_\rho(f)-S_\rho(g)=-2D+\rho q+\mathcal R,\qquad
 0\leq\mathcal R\leq C_Q\rho D,\qquad C_Q=C_dM_+.
\end{equation}
This bounds the perturbation's quadratic remainder; it does not estimate
the difference of the two full higher-chaos tails separately.

Here the original mass error $u$ in the repair lemma is zero. Thus
\eqref{eq:repair-bounds}--\eqref{eq:repair-coercivity} yield
\[
 \rho q\leq(2-\sigma)\rho J_{\rm in}+E_RD,\qquad
 E_R=2C_J\eps_R+36(2K+1)A_+^2\eps_R^2.
\]
Choose a finite $R\geq3$ with $E_R\leq\sigma/4$, possible because
$\eps_R\to0$. Choose any positive $\rho_m^+$ satisfying
\begin{equation}\label{eq:positive-range}
 \rho_m^+\leq
 \min\left\{\frac1{16},\frac{\sigma}{16\beta_R},
                  \frac{\sigma}{4C_Q},\frac{c_A^2}{576L}\right\}.
\end{equation}
For every competitive assignment at $0<\rho\leq\rho_m^+$,
$\eta\leq c_A/24$, so the repair is available.
Using~\eqref{eq:compact} and $(1-\beta_R\rho)^{-1}\leq1+2\beta_R\rho$,
\[
 \rho q\leq\left(\frac{2-\sigma}{1-\beta_R\rho}+E_R\right)D.
\]
Since $4\beta_R\rho\leq\sigma/4$, substituting this into
\eqref{eq:positive-comparison} proves
\[
 S_\rho(f)-S_\rho(g)\leq-\frac{\sigma}{4}D.
\]
Strict ordering in~\eqref{eq:plus-gap} makes $D>0$ unless $f=g$.
Thus no distinct competitive assignment exists in dimension $d$.
The finite-dimensional and equality transfer is completed in
Section~\ref{sec:dimension}. The interval is closed at its upper endpoint
because every defining bound in~\eqref{eq:positive-range} is non-strict.


\section{The negative transfer with free common masses}\label{app:negative}
We first work in dimension $d=m-1$. A pair is competitive if
$B_{-r}(f,h)\leq B_{-r}(g,g)$, and its matching mass vector is denoted by $p$.
Let $F=(f+h)/2$ and $U=\1/m$.

\subsection{Uniform localization in the nearest frame}
Parseval gives
$\norm{b(f)}_{\HS}^2\leq\sum_i\Var(f_i)\leq1-\norm p^2\leq1$,
and likewise for $h$.
The absolute value of the degree-at-least-two bilinear Hermite remainder
is at most $r^2$, by Cauchy--Schwarz; the same bound holds for the reference.
Competitiveness therefore implies
\begin{equation}\label{eq:competitive-negative}
 \norm{p-U}^2\leq r(\ip{b(f)}{b(h)}_{\HS}-w)+2r^2\leq r+2r^2,
 \qquad
 \ip{b(f)}{b(h)}_{\HS}\geq w-2r.
\end{equation}
In particular all masses are positive if $r\leq1/(4m^2)$.

For every $\eta_*>0$, there is a $\delta_*>0$ such that all competitive
pairs with $0<r<\delta_*$ obey
\begin{equation}\label{eq:qual-local}
 \frac12\sum_i\norm{F_i-g_i}_1<\eta_*,
\end{equation}
where $g$ is selected by \emph{any} nearest regular centroid matrix to $b(F)$.
To prove this uniformly, suppose otherwise and choose a violating sequence
$r_n\downarrow0$. Weak compactness of the bounded simplex-valued set in
$L^2(\gamma_d)^m$ gives weak subsequential limits $f_*,h_*$.
By~\eqref{eq:competitive-negative}, they are balanced and have centroid
inner product at least $w$. The endpoint and Cauchy--Schwarz give the reverse
chain
\[
 \ip{b(f_*)}{b(h_*)}\leq\norm{b(f_*)}\norm{b(h_*)}\leq w.
\]
Equality forces identical centroid matrices of norm squared $w$, and endpoint
equality gives $f_*=h_*=g_*$ for one ordered regular partition.
Weak convergence to an indicator gives strong $L^1$ convergence, since
$\sum_i\norm{f_{n,i}-g_{*,i}}_1=2-2\sum_i\ip{f_{n,i}}{g_{*,i}}\to0$,
and likewise for $h_n,F_n$.
The centroids converge by weak continuity of first moments.
Any nearest regular matrices $c_n$ then satisfy
$\norm{c_n-b(F_n)}\leq\norm{c_*-b(F_n)}\to0$.
Their winning labels converge off the null hyperplanes of $c_*$, so their
regular partitions converge in $L^1$. This contradicts the violation and
proves~\eqref{eq:qual-local}.

\subsection{The mixed comparison retains the mass penalty}
Use a nearest regular frame for $b(F)$, and put
\[
 D=D_-(F),\quad z=\sqrt{D/r},\quad
 \eta=\tfrac12\sum_i\norm{F_i-g_i}_1,\quad
 q=\norm{b(F)-c}_{\HS}^2,\quad u=\tfrac12\norm{p-U}_1.
\]
The loss $D_-$ is affine, so $2D=D_-(f)+D_-(h)$.
Writing $q_f=f-g$ and $q_h=h-g$, the constant term in their mixed quadratic
expansion is $\norm{p-U}^2$. Its first-order term satisfies
\[
 \ip{b(f)-c}{b(h)-c}
 =q-\tfrac14\norm{b(f)-b(h)}_{\HS}^2\leq q.
\]
The higher remainder has absolute value at most
$C_Qr\sqrt{D_-(f)D_-(h)}\leq C_QrD$, where $C_Q=C_dM_-$.
Consequently
\begin{equation}\label{eq:negative-comparison}
 B_{-r}(f,h)-B_{-r}(g,g)
 \geq\norm{p-U}^2+2D-rq-C_QrD.
\end{equation}
No positivity of a negative-correlation higher-chaos term is assumed.

\subsection{Repair and absorption of the unequal masses}
Choose the localization cutoff~\eqref{eq:qual-local} with $\eta_*=c_A/24$.
The repair in Appendix~\ref{subsec:repair} is then available. Its mass
imbalance satisfies both $v\leq u+\eps_RA_-z$ and $v\leq\eta$.
Using~\eqref{eq:repair-bounds} and~\eqref{eq:repair-coercivity},
$\eta\leq A_-z$, and $(x+y)^2\leq2x^2+2y^2$, we obtain
\begin{align}
 rq&\leq(2-\sigma)rJ_{\rm in}+B\sqrt{rD}\,u
                  +72(2K+1)ru^2+\widetilde E_RD,\label{eq:negative-rq}\\
 B&=6A_-/c_A,\qquad
 \widetilde E_R=2C_J\eps_R+72(2K+1)A_-^2\eps_R^2.\nonumber
\end{align}
Young's inequality gives
$B\sqrt{rD}\,u\leq\sigma D/8+(2B^2/\sigma)ru^2$.
Put $C_{\rm mass}=72(2K+1)+2B^2/\sigma$.

Choose finite $R\geq3$ with $\widetilde E_R\leq\sigma/8$.
Then choose $r_m^->0$ small enough that
\begin{equation}\label{eq:negative-range}
 r_m^-\leq\min\left\{\frac1{16},\frac1{4m^2},
 \frac{\delta_*}{2},\frac{\sigma}{32\beta_R},
 \frac{\sigma}{8C_Q},\frac{2}{mC_{\rm mass}}\right\}.
\end{equation}
Here $\delta_*$ is the uniform cutoff just proved. Every entry is positive
and depends only on $m$ and its finite coercivity constant.
For $0<r\leq r_m^-$, the compact comparison~\eqref{eq:compact} and
$(1-\beta_Rr)^{-1}\leq1+2\beta_Rr$ imply from~\eqref{eq:negative-rq}
\[
 rq\leq(2-5\sigma/8)D+C_{\rm mass}ru^2.
\]
Insert this into~\eqref{eq:negative-comparison} and use
$u^2\leq m\norm{p-U}^2/4$ to get
\begin{equation}\label{eq:negative-strict}
 B_{-r}(f,h)-B_{-r}(g,g)
 \geq\tfrac12\norm{p-U}^2+\tfrac{\sigma}{2}D.
\end{equation}
Thus a competitive pair has $p=U$ and $D=0$.
The strict gap~\eqref{eq:minus-gap} forces $F=g$; since $g$ is an indicator
and $F=(f+h)/2$, simplex-valuedness forces $f=h=g$ almost everywhere.
This proves the full relaxed fixed-dimensional assertion. The remaining
dimension reduction and recovery of fractional equality are given next in
Section~\ref{sec:dimension}. No quantitative localization for unbalanced
assignments was assumed: its uniform qualitative form was proved above.


\section{Proof lineage and review boundaries}\label{app:lineage}
The manuscript contains the covariance and endpoint proofs in full.
Appendices~\ref{app:common}--\ref{app:negative} also contain the fixed-dimensional
transfer proofs, including their constants and closed-endpoint choices.
The analytic external inputs are the convex-support variance inequality
\cite[Proposition 2.1]{BL2000}, the two dimension reductions with distinct mass
scopes stated in Section~\ref{sec:dimension}, and, for the discrete corollary,
the same-correlation transfer and plurality limit in \cite{IM2009}.
No computational covariance cover is a premise of these results.

The assembly uses the following local sources, preserved separately in the
review project's proof-source and audit directories. Hashes are SHA256; the
companion reproduction manifest records the complete file inventory.
\begin{enumerate}
\item The all-dimensional covariance proof:
\path{ALL_DIMENSIONAL_COVARIANCE_CONCAVITY.txt}.
\begin{quote}\small\ttfamily
225db53f4489bf513926eead05cc5ec8\\
c40962991246ac5036a897faa2d8130b
\end{quote}
\item The general positive transfer:
\path{FINITE_CELL_COERCIVITY_TO_NOISE_TRANSFER_20261030.txt}.
\begin{quote}\small\ttfamily
1d951abe62965c31d3b8ce5209099d62\\
167ed9fc9b554f0c91d4c057b9a3dff5
\end{quote}
\item The general free-common-mass negative transfer:
\path{FREE_COMMON_MASS_NEGATIVE_TRANSFER_AUDIT_20261031.txt}.
\begin{quote}\small\ttfamily
9a6339628de87b4d0ae95526438caab8544\\
da303bb2934c8494c204b33adc2ef
\end{quote}
\end{enumerate}
The older transfer notes have conditional historical status wording.
Theorem~\ref{thm:endpoint} supplies precisely their hypothesis
\eqref{eq:Hm} for every fixed $m$; it is no longer an unproved endpoint
assumption in the arguments presented here.

The accompanying records include the full cleanroom covariance audit,
the separate end-to-end application audit, the exact scope note, and the
root implication note. These are internal AI-agent audits; no external
reviewer is implied. The independent geometry checker verifies the
facet-basis construction, the full chain factorization, and the redundant
gauge equations on rational fixtures in dimensions two through seven.
Its 96 exact checks support the algebra but do not replace the uniform
proof or the analytic hypotheses. Earlier frozen four-cell releases retain
their original statements and numerical intervals; they are not dependencies
of the new all-cell analytic covariance proof.

This project is a local review artifact. Internal AI-agent mathematical audits
and reproducibility checks are distinguished from external peer review.
No author or affiliation has been assigned, and no publication, external
contact, or priority claim is made. The main conclusions are for each fixed
cell count near independence; full positive and negative correlation
ranges remain open.


\begin{thebibliography}{9}
\bibitem{H2014}
S. Heilman,
\emph{Euclidean Partitions Optimizing Noise Stability},
2014.
\href{https://arxiv.org/abs/1211.7138v2}{arXiv:1211.7138v2}.
Section 8, Conjecture 3 and the preceding discussion.

\bibitem{H2026}
S. Heilman,
\emph{On the Propeller Conjecture},
23 October 2026.
\href{https://arxiv.org/abs/2610.32158v1}{arXiv:2610.32158v1}.
Introduction following Theorem 1.3.

\bibitem{BL2000}
S.~G. Bobkov and M. Ledoux,
\emph{From Brunn--Minkowski to Brascamp--Lieb and to logarithmic Sobolev
inequalities}, Geometric and Functional Analysis \textbf{10} (2000),
1028--1052.
Proposition 2.1, printed p.~1031.
\href{https://www.math.univ-toulouse.fr/~ledoux/BrunnMinkowski.pdf}{Author-hosted full text}.

\bibitem{HT2021}
S. Heilman and A. Tarter,
\emph{Three Candidate Plurality is Stablest for Small Correlations},
Forum of Mathematics, Sigma \textbf{9} (2021), e65.
\href{https://arxiv.org/abs/2011.05583v2}{arXiv:2011.05583v2}.
Theorem 1.9.

\bibitem{H2022}
S. Heilman,
\emph{Hyperstable Sets with Voting and Algorithmic Hardness Applications},
2022.
\href{https://arxiv.org/abs/2209.11216v1}{arXiv:2209.11216v1}.
Problem 7.1 and Theorem 7.9.

\bibitem{IM2009}
M. Isaksson and E. Mossel,
\emph{Maximally Stable Gaussian Partitions with Discrete Applications},
2009.
\href{https://arxiv.org/abs/0903.3362v3}{arXiv:0903.3362v3}.
Theorem 7.1 and Lemma 7.3.
\end{thebibliography}

\end{document}
